
\section{Census as a Runnable Tool: \texttt{dx-census}}
\label{sec:dxcensus}
The annotation-consistency audit is not a one-off analysis in this paper --- it is a
command in the shipped tool, with the admissibility gate enforced in code:
\begin{verbatim}
usage: dx-census [-h] [--image-size IMAGE_SIZE] [--epochs EPOCHS]
                 [--folds FOLDS] [--seed SEED]
                 folder
\end{verbatim}
Given a folder of class-named images, \texttt{dx-census} trains the
out-of-fold probe, runs the confident-learning estimator, applies the gate
and \textbf{refuses to emit a census} when the gate fails, printing the
diagnostic instead. This behaviour was demonstrated on the inadmissible
\texttt{tissuemnist} probe (the external-battery section): the tool reported the
fabrication risk rather than a noise number. The gate thresholds
($\mathrm{OOF\ acc}\ge\max(0.60,\ \mathrm{majority}+0.05)$, $\ge 250$
optimizer steps per fold) are constants in the source, documented, and
covered by a unit test using the pneumonia gate values.
The census output is a JSON with per-image flags (exact paths), the
confusion-matrix estimate of joint label distribution, per-direction
contamination rates, and the gate verdict. Every number in
the census section of this paper was produced by this command on the
public archive, not by bespoke analysis code.

\section{What Would Falsify Our Claims}
\label{sec:falsify}
We state the falsifiability conditions explicitly, as a matter of method:
\begin{enumerate}
\item \textbf{Census rate.} An independent re-annotation of our 155
published flags by blinded expert radiologists showing $<50\%$
agreement would refute the noise estimate's face validity; the flag list
is published precisely to enable this.
\item \textbf{Directionality.} A re-estimation with a different
out-of-fold model family (e.g., a pretrained transformer probe) that
recovers a symmetric confusion structure would weaken the one-directional
contamination claim; we report our probe's architecture and the
crosscheck model's agreement ($J=0.841$ on the malaria census) as the
current support.
\item \textbf{Report-parsing signature.} The hypothesis that label
contamination co-varies with acquisition quality predicts that flagged
films are blurrier and flatter (verified here, $p<10^{-6}$). A larger
acquisition-metadata study showing flagged films are \emph{not}
acquisition-degraded would refute the mechanism, though not the census
rate.
\item \textbf{Cleaning gain.} The +5.9-point gain was measured on one
seed and one architecture pair; a multi-seed replication showing the
gain within noise would demote it from finding to anecdote. We flag this
as the highest-priority follow-up experiment.
\end{enumerate}

\section{Discussion}
The three findings of this paper form a coherent account of
\texttt{ChestXRay2017}. Its labels were generated by automated parsing of
radiology reports; report text understates or omits findings that are
visually present on the image, producing NORMAL labels on films a model
reads as PNEUMONIA --- the 42:1 asymmetry we measure. The same borderline
films that strain report parsing are disproportionately the blurred,
low-dynamic-range acquisitions (the property stratification), which is
why the contamination is not a random sample of the dataset but a
structured, measurable stratum. Cleaning the stratum improves a compact
model by +5.9 points because the removed labels were actively misleading
gradients; on malaria, where the expert-audited test set shares no
contamination with training, the same procedure does not transfer.
The broader implication for medical-imaging benchmarks: a published
accuracy on a web-mined or report-parsed dataset is a property of the
label process as much as of the model, and an annotation-consistency audit with an admissibility gate is a cheap prerequisite ($\sim$2 CPU-hours here) for
interpreting any such number. We recommend benchmark papers report a
census alongside the leaderboard, and we ship the tool that produces one.
